\newtheorem{definicion}{Definición}
\newtheorem{ejemplo}{Ejemplo}
\documentclass{beamer}
\usepackage[spanish]{babel}
\usepackage[utf8]{inputenc}
\usepackage{graphicx} 

\title[Presentacion con Beamer]{Series de pontencias de Taylor: ln(x)} %No lo ponemos dentro del begin porque ya están predefinidos
\author[Técnicas experimentales{ULL}]{Yoselin Armas Ramos y Bianca E. Kennedy Giménez}
\date[16-05-14]{16 de Mayo de 2014}

\usetheme{Madrid}

\begin{document}

\begin{frame}

  \includegraphics[width=0.15\textwidth]{img/logotipo-secundario-ULL}
  \hspace*{7.0cm}
  \includegraphics[width=0.16\textwidth]{img/fmatesc}
  \titlepage
  
  \begin{small}
    \begin{center}
     Facultad de Matemáticas \\
     Universidad de La Laguna
    \end{center}
  \end{small}
\end{frame}


\begin{frame}
\frametitle{Índice}
\tableofcontents[pausesections] %con este comando hacemos que ponga en esta transparencia todas las secciones [opcion para que aparezcan con pausa en casa clic]
\end{frame}

\begin{frame}
\section{Historia} %Abrimos una sección
\frametitle{Historia}
\begin{itemize}
\item El filósofo griego Zenón y sus paradojas.
\item Revolución filosófica.
\item Siglo VXI,Madhava de Sangamagrama y los primeros usos de la serie de Taylor.
\item James Gregory y su interés por las Series de Taylor,series de Maclaurin.
\item En el siglo XVIII ,el matemático británico Brook Taylor presentó de manera formal el Desarrollo de Taylor.
\end{itemize}
\end{frame}

\begin{frame}
\subsection{Series de Taylor}
\frametitle{Series de Taylor.}
\begin{figure}[!th]
\begin{center}
\includegraphics[width=0.20\textwidth]{images/taylor.eps}
\caption{Brook Taylor}
\label{fig:1}
\end{center}
\end{figure}
\begin{itemize}
\item El desarrollo de Taylor parte de una serie infinita y da como resultado un resultado finito.
\item Polinomio de Taylor:
\end{itemize}
 \[ P_{(n,a)} = f(a)+\frac{f'(a)}{1!}(x-a)+\frac{f''(a)}{2!}(x-a)^2+\cdots+\frac{f^{(n)}(a)}{n!} (x-a)^{n}\,. \]
\end{frame}


\begin{frame}
\frametitle{Logaritmo Neperiano} %No es posible sin abrir una sección
Estudiando los fenómenos de creciemiento y decrecimiento en la naturaleza, se observó que con frecuencia aprecían potencias de un número irracional al que se llamo ``e'', 
cuyo valor aproximado es:
 \[ e \approx 2,7182818284590452353602874713527. \] 

Para estudiar estos fenómenos se aplicaban los logaritmos y sus propiedades, en concreto el logaritmo en base ``e'', también llamado ``logaritmo neperiano''. 
\[ln(x)=log_e(x)\]
\end{frame}



\section{Aplicación de la serie de Taylor de Ln(x)}

\begin{frame}
\frametitle{Aplicación de la serie de Taylor de Ln(x)}
Se aplica el desarrollo de Taylor de manera general, donde:
\[n=grado, a=valor, f^t(a)= \frac{b}{a^t} (1<t<n)\]

Sea f(x)=ln(x) :

\begin{enumerate}
\item Se calcula la imagen:
f(a) = ln(a)
\item Se calcula la primera derivada y su imagen:
$f'(a) = \frac{1}{a}$ 
\item Se calcula la segunda derivada y su imagen:
$f''(a) = \frac{-1}{a^2}$
\item Se calcula la tercera derivada y su imagen:
$f'''(a) = \frac{2}{a^3}$
\item Calculamos la enésima derivada y su imagen:
$f^n(a) = \frac{(n-1).(-b)}{a^{(n)}}\,.$
\end{enumerate}
Entonces,
\[ P_{(n,a)} = ln(a)+\frac{\frac{1}{a}}{1!}(x-a)+\frac{\frac{-1}{a^2}}{2!}(x-a)^2+\cdots +\frac{ \frac{(n-1).(-b)}{a^{(n)}}}{n!} (x-a)^{n}\,. \]
\end{frame}

\section{Proceso experimental}
\subsection{Variando el centro}
\begin{frame}
A continuación, veremos como variando solo el centro el desarrollo de Taylor cambia sus valores notablemente

\begin{tabular}{|l | r@{,}l |}
\hline
Si el valor de c= 2 &  1 & 69314718055995\\
\hline
Si el valor de c=4  & -7746 & 38037230555\\
\hline
Si el valor de c=6  & -1680074 & 89712942
\\
\hline
\end{tabular}

\end{frame}

\subsection{Variando el grado}
\begin{frame}

PONER TABLA CON VARIACIONES EN EL GRADO

\end{frame}


\section{Apéndice}

\subsection{Python}
\begin{frame}[fragile]
A continuación,mostramos un programa elaborado en Python,un lenguaje interpretado orientado a objetos.Este programa realiza las series de Taylor de grado n centrado en un punto c de la función Ln(x).
Para que el programa nos devuelva el resultado debemos darle los valores deseados al punto,grado y al centro. 
\begin{verbatim}
#!/usr/bin/python
#!encoding: UTF-8

from math import*
from sympy import*     

c = symbol('c')   
funcion=ln(c)

#Declaramos la función factorial.
def fac(n):    
    if n == 0:
      return 1
    else:
      return n * fac(n-1)

#Declaramos la serie de Taylor.
def Taylor(grado, punto, centro, funcion):   
  suma = funcion.evalf(subs={c:centro}) 
  
#Hacemos un bucle for para el cálculo de derivadas. 
  for i in range(1,grado+1):                
    deriv = diff(funcion, c)
    termino = (deriv.evalf(subs={c:centro})/fac(i))*(punto-centro)**i
    funcion = deriv
    suma += termino
  return suma  
\end{verbatim}
\end{frame}


\subsection{Información de la máquina}
\begin{frame}[fragile]
A continuación, mostramos la información del ordenador donde se realizaron los experimentos:

\begin{verbatim}
('default', 'Sep 26 2013 20:03:06')
<function platform at 0x7fe3b32c15f0>
('Linux', 'bianca-VPCEA2S1E', '3.8.0-34-generic',
'#49~precise1-Ubuntu SMP Wed Nov 13 18:05:00 UTC 2013', 'x86_64', 'x86_64')
2.7.3
Intel(R) Core(TM) i3 CPU   M 350  @ 2.27GHz
GenuineIntel
2266.000 Hz 
3072 KB
\end{verbatim}
\end{frame}

\begin{frame}
\frametitle{Bibliografía}
\begin{thebibliography}
\beamertemplatebookbibitems
\bibitem[Guía Docente,2013]{guia}
Guía docente (año 2013)
{\small $http://gjtsrh.com$}
\end{thebibliography}
\end{frame}
\end{document}

